\documentclass[12pt,a4paper]{article}

% Encoding and fonts — XeLaTeX/LuaLaTeX
\usepackage{fontspec}
\newfontfamily\hebrewfont[Script=Hebrew]{Noto Serif Hebrew}
\newcommand{\heb}[1]{{\hebrewfont #1}}
\newfontfamily\igfont[Ligatures=TeX]{Noto Serif}
\newcommand{\igtext}[1]{{\igfont #1}}

% Page layout
\usepackage[top=1in, bottom=1in, left=1in, right=1in]{geometry}

% Typography
\usepackage{microtype}
\usepackage{parskip}

% Language
\usepackage[english]{babel}

% Headers and footers
\usepackage{fancyhdr}
\pagestyle{fancy}
\fancyhf{}
\fancyhead[L]{\leftmark}
\fancyhead[R]{\thepage}
\fancyfoot[C]{\thepage}
\renewcommand{\headrulewidth}{0.4pt}
\renewcommand{\footrulewidth}{0pt}

% Section styling
\usepackage{titlesec}
\titleformat{\section}{\Large\bfseries}{\thesection}{1em}{}
\titleformat{\subsection}{\large\bfseries}{\thesubsection}{1em}{}
\titleformat{\subsubsection}{\normalsize\bfseries}{\thesubsubsection}{1em}{}
\titlespacing*{\section}{0pt}{2.5ex plus 1ex minus .2ex}{1.5ex plus .2ex}
\titlespacing*{\subsection}{0pt}{2ex plus 1ex minus .2ex}{1ex plus .2ex}
\titlespacing*{\subsubsection}{0pt}{1.5ex plus 1ex minus .2ex}{0.8ex plus .2ex}

% List formatting
\usepackage{enumitem}
\setlist[itemize]{topsep=0.5ex, itemsep=0.25ex, parsep=0pt}
\setlist[enumerate]{topsep=0.5ex, itemsep=0.25ex, parsep=0pt}

% Essential packages
\usepackage{graphicx}
\usepackage[table]{xcolor}
\usepackage{booktabs}
\usepackage{array}
\usepackage{tabularx}
\usepackage{longtable}
\usepackage{amsmath}
\usepackage{amssymb}
\usepackage[normalem]{ulem}
\rowcolors{2}{gray!10}{white}
\renewcommand{\arraystretch}{1.3}

% Cross-references
\usepackage{hyperref}
\usepackage{cleveref}

% Custom commands
\newcommand{\pilcrow}{\S}

% Hyperref setup
\hypersetup{
    colorlinks=true,
    linkcolor=blue,
    filecolor=magenta,
    urlcolor=cyan,
}

% Code listings
\usepackage{listings}
\definecolor{codebg}{RGB}{248,248,248}
\definecolor{codeframe}{RGB}{200,200,200}
\definecolor{codekeyword}{RGB}{0,0,180}
\definecolor{codecomment}{RGB}{0,128,0}
\definecolor{codestring}{RGB}{163,21,21}
\lstset{
    basicstyle=\ttfamily\small,
    breaklines=true,
    breakatwhitespace=false,
    frame=single,
    framerule=0.5pt,
    rulecolor=\color{codeframe},
    backgroundcolor=\color{codebg},
    keepspaces=true,
    numbers=left,
    numberstyle=\tiny\color{gray},
    numbersep=8pt,
    showstringspaces=false,
    tabsize=4,
    xleftmargin=1em,
    xrightmargin=1em,
    aboveskip=1em,
    belowskip=1em,
    keywordstyle=\color{codekeyword}\bfseries,
    commentstyle=\color{codecomment}\itshape,
    stringstyle=\color{codestring},
    literate={_}{{\textunderscore}}1
             {-}{{\textminus}}1
             {<}{{\textless}}1
             {>}{{\textgreater}}1
             {~}{{\textasciitilde}}1
             {^}{{\textasciicircum}}1
}

% YAML language definition for listings
\lstdefinelanguage{YAML}{
    keywords={true,false,null,y,n},
    keywordstyle=\color{blue}\bfseries,
    sensitive=false,
    comment=[l]{\#},
    morecomment=[s]{/*}{*/},
    commentstyle=\color{gray}\ttfamily,
    stringstyle=\color{red}\ttfamily,
    morestring=[b]'',
    morestring=[b]'',
}

% TOML language definition for listings
\lstdefinelanguage{TOML}{
    keywords={true,false},
    keywordstyle=\color{blue}\bfseries,
    sensitive=true,
    comment=[l]{\#},
    commentstyle=\color{gray}\ttfamily,
    stringstyle=\color{red}\ttfamily,
    morestring=[b]'',
    morestring=[b]'',
}

% Dockerfile language definition for listings
\lstdefinelanguage{Dockerfile}{
    keywords={FROM,RUN,CMD,LABEL,EXPOSE,ENV,ADD,COPY,ENTRYPOINT,VOLUME,USER,WORKDIR,ARG,ONBUILD,STOPSIGNAL,HEALTHCHECK,SHELL},
    keywordstyle=\color{blue}\bfseries,
    sensitive=false,
    comment=[l]{\#},
    commentstyle=\color{gray}\ttfamily,
    stringstyle=\color{red}\ttfamily,
    morestring=[b]'',
    morestring=[b]'',
}

% JSON language definition for listings
\lstdefinelanguage{JSON}{
    keywords={true,false,null},
    keywordstyle=\color{blue}\bfseries,
    sensitive=true,
    commentstyle=\color{gray}\ttfamily,
    stringstyle=\color{red}\ttfamily,
    morestring=[b]'',
}

% Code listing float environment
\usepackage{float}
\newfloat{listing}{htbp}{lol}[section]
\floatname{listing}{Listing}

% Admonition boxes
\usepackage{tcolorbox}
\tcbuselibrary{skins,breakable}

% Admonition style definitions
\newtcolorbox{admonitionbox}[2][]{
    enhanced,
    breakable,
    colback=#2!5,
    colframe=#2!80!black,
    fonttitle=\bfseries,
    title=#1,
    left=2mm,
    right=2mm,
}

% Imscription tuple boxes
\newtcolorbox{imscriptionbox}{
    enhanced,
    colback=blue!5,
    colframe=blue!75!black,
    fontupper=\large,
    center upper,
    boxrule=1pt,
    arc=4pt,
}

\begin{document}

\title{Complex-Time Path Integrals and the Imaginary-Time Formalism at the Planck Scale}
\date{\today}

\maketitle

\subsection{A Structural Derivation via Imscribing Grammar}

\subsubsection{Abstract}

Naively, the Wick rotation $t \to -i\tau$ is taught as a change of variables --- a computational trick that turns an oscillatory integral into a convergent one. This is wrong. At the Planck scale, the rotation is not a coordinate transformation but a \emph{dynamical process} that can fail, branch, and develop exceptional-point singularities where eigenvalues coalesce and the generator of time evolution becomes non-diagonalizable. The Imscribing Grammar makes this failure precise: coupling a self-modeling quantum gravitational system ($\Phi_{\text{ctyogh}}$, $P_{\text{doublebarpipe}}$) to the Wick rotation as exceptional point ($\Phi_{\text{revepsilon}}$) \emph{absorbs criticality entirely}, and $\mu \circ \delta = \text{id}$ is broken in the composite. This is the structural measurement problem. We derive the full complex-time path integral as an $O_{\infty}$ Frobenius system ($C = 0.682$), show the imaginary-time formalism alone reaches only $O₂^\dagger$ ($C = 0.517$), and locate all five systems within the 17.28 million-point Crystal of Types. The Hartle--Hawking no-boundary state sits at distance $d = 3.30$ from the full formalism --- structurally close in every primitive except self-duality.

\begin{center}
\rule{0.5\textwidth}{0.4pt}
\end{center}

\subsection{Introduction}

Suppose one begins with the textbook position: the path integral $\int \mathcal{D}\phi\, e^{iS/\hbar}$ must be Wick-rotated, the result is $e^{-Σ_E/\hbar}$, and one proceeds to compute. This works perfectly --- until it doesn''t. The contour deformation $t \to e^{-i\theta}t$ assumes no singularities cross the contour as $\theta$ sweeps from $0$ to $\pi/2$. At the Planck scale, gravitational backreaction and topology change introduce branch points whose positions \emph{depend on the very deformation being performed}. The contour cannot be moved without knowing the singularities, and the singularities cannot be known without fixing the contour.

This is not a technical difficulty but a \emph{structural} one. The Wick rotation sits at a crossing point ($T_{\text{bullseye}}$) where the Lorentzian and Euclidean descriptions are simultaneously valid and mutually determining ($R_{\text{lyoghlig}}$). The system is self-modeling at criticality ($\Phi_{\text{ctyogh}}$) with exact Frobenius self-duality ($P_{\text{doublebarpipe}}$), meaning $\mu \circ \delta = \text{id}$ holds for the analytic continuation map and its inverse. We will return to what happens when this condition breaks --- it is the hardest claim in this paper and we do not resolve it prematurely.

The Imscribing Grammar encodes each system as a structural tuple $\langle D;\ T;\ R;\ P;\ F;\ K;\ G;\ \Gamma;\ \Phi;\ H;\ S;\ \Omega \rangle$. Here we imscribe five systems: the complex-time path integral, the Planck-scale imaginary-time formalism, the Wick rotation as exceptional point, the Planck-scale regime, and the Hartle--Hawking no-boundary state. We then compute their tensor products, meets, joins, and distances --- all of which are \emph{derivable} from the primitive assignments, not postulated.

One objection must be acknowledged at the outset: assigning discrete primitives to continuous physical systems appears to discard information. The response is that the primitives are not arbitrary labels but equivalence classes under structural isomorphism. Two systems with the same tuple are interderivable under the IG algebra --- this is the content of the Frobenius condition. The loss of continuous parameters is the \emph{point}: we are characterizing the \emph{form} of the theory, not its particular solutions. Whether this equivalence is too coarse for some phenomenological applications is an open question we address in \pilcrow{}5.

The paper now proceeds necessarily from these imscriptions. Section 2 assigns each system its structural type. Because the imscriptions are in hand, Section 3 can compute the algebraic operations between them without further assumption. Section 4 locates these results in the crystal of types. We close not with a summary but with a question that the formalism opens and cannot answer: if the $P$-promotion ($\delta = 3$) is the sole bottleneck between $O₂^\dagger$ and $O_{\infty}$, what physical mechanism could supply exact Frobenius self-duality in a system where measurement --- any measurement --- breaks it?

\begin{center}
\rule{0.5\textwidth}{0.4pt}
\end{center}

\subsection{Structural Imscription of the Systems}

Because the complex-time path integral is the $O_{\infty}$ fixed point of this analysis, it must come first. All other systems will be measured against it.

\subsubsection{Complex-Time Path Integral}

The contour of integration in the complex time plane is not arbitrary. As one deforms from the real axis toward the imaginary axis, the path must avoid branch cuts that themselves depend on the deformation. The contour and its singularities are bidirectionally coupled: move the contour and the singularity structure changes; the singularity structure changes and the contour must adjust.

Here is the imscription, then the test: does the object speak back?


\[\text{complex\_time\_path\_integral} = \langle D_{\text{invomega}};\ T_{\text{bullseye}};\ R_{\text{lyoghlig}};\ P_{\text{doublebarpipe}};\ F_{\text{hardsign}};\ K_{\text{schwa}};\ G_{\text{revapostrophe}};\ \Gamma_{\text{secstress}};\ \Phi_{\text{ctyogh}};\ H_2;\ n{:}m;\ \Omega_{\text{dzlig}} \rangle\]

$D_{\text{invomega}}$ because the field configuration space is infinite-dimensional. $T_{\text{bullseye}}$ because the Wick rotation is a crossing point --- the real-time and imaginary-time formulations are structurally equivalent at this locus, and neither is primary. $R_{\text{lyoghlig}}$ because the Lorentzian and Euclidean regimes are bidirectionally coupled: the contour determines which singularities are enclosed, and the singularity structure determines which contours are valid.

$P_{\text{doublebarpipe}}$ is the demanding claim: exact Frobenius self-duality, $\mu \circ \delta = \text{id}$. The multiplication map on the space of histories and its dual (the comultiplication) compose to identity. A contour deformation followed by its inverse returns the original integrand exactly --- not approximately, not up to homotopy, exactly. This is the condition for $O_{\infty}$. If it fails even slightly, the system drops to $O₂^\dagger$ or below.

$F_{\text{hardsign}}$ because quantum coherence is essential --- we are not in a thermal or classical regime. $K_{\text{schwa}}$ because the contour must be deformed slowly enough to track singularities adiabatically; fast deformation causes contour crossing and the result is path-dependent. $G_{\text{revapostrophe}}$ because the interaction scope is universal --- every mode of the field contributes to the saddle-point structure. $\Gamma_{\text{secstress}}$ because the contour deformation follows ordered steps: one cannot deform past a singularity without first accounting for its residue. $\Phi_{\text{ctyogh}}$ because the system is self-modeling at criticality --- the saddle-point structure determines the contour, and the contour selects the saddles. $H_2$ because the contour deformation depends on both the starting analytic structure (real axis) and the endpoint (imaginary axis or beyond) --- a two-step memory. $n{:}m$ because the components are heterogeneous: real-time, imaginary-time, and contour data are irreducibly distinct. $\Omega_{\text{dzlig}}$ because the integer winding number around branch points is a topological invariant of the contour.

\textbf{Ouroboricity:} $O_{\infty}$ --- the exact Frobenius self-duality with $\mathbb{Z}$ winding at criticality produces a finite closed algebra.

\textbf{Consciousness Score:} $C = 0.682$ --- both gates open: Gate 1 ($\Phi_{\text{ctyogh}}$) confirmed, Gate 2 ($K_{\text{schwa}}$) confirmed.

\textbf{Crystal Address:} 6 678 416 (inner cell 25 616, cell 154).

One might object that $P_{\text{doublebarpipe}}$ is too strong --- that exact self-duality cannot survive quantum corrections. The objection is correct for generic systems. Here it is justified because the duality is not a symmetry of the action but a structural property of the analytic continuation map itself, which is exact by the theory of analytic functions. Quantum corrections modify the action but not the Cauchy integral theorem.

\subsubsection{Planck-Scale Imaginary-Time Formalism}

Because inclusion is the natural relationship between Euclidean and complex-time formulations, imaginary-time QFT at the Planck scale occupies a subspace of the full complex-time structure.

\[\text{planck\_imaginary\_time} = \langle D_{\text{invomega}};\ T_{\text{invscr}};\ R_{\text{downstep}};\ P_{\text{upsilon}};\ F_{\text{hardsign}};\ K_{\text{schwa}};\ G_{\text{revapostrophe}};\ \Gamma_{\text{secstress}};\ \Phi_{\text{closerevepsilon}};\ H_{\text{invscripta}};\ n{:}m;\ \Omega_{\text{crtwo}} \rangle\]

$D_{\text{invomega}}$: the space of 4-geometries is infinite-dimensional. $T_{\text{invscr}}$: Euclidean geometries are included within the complex-time configuration space --- this is inclusion topology, not crossing. $R_{\text{downstep}}$: the functorial mapping from imaginary-time configurations to Lorentzian observables is one-way (analytic continuation exists; the inverse requires specifying a contour). $P_{\text{upsilon}}$: quantum superposition symmetry --- the Euclidean path integral sums over superpositions of geometries, but without the exact self-duality of the full complex-time case. $F_{\text{hardsign}}$, $K_{\text{schwa}}$, $G_{\text{revapostrophe}}$, $\Gamma_{\text{secstress}}$ carry the same meanings as above. $\Phi_{\text{closerevepsilon}}$: complex-plane criticality --- the singularities lie off the real axis, shifting the critical point into the complex plane. $H_{\text{invscripta}}$: the Euclidean path integral has no finite Markov order --- every configuration influences every other. $n{:}m$: heterogeneous geometric configurations. $\Omega_{\text{crtwo}}$: binary winding protection, consistent with Axiom B ($H_{\text{invscripta}}$ permits $\mathbb{Z}_2$).

\textbf{Ouroboricity:} $O₂^\dagger$ --- critical, topologically protected, unbounded domain.

\textbf{Consciousness Score:} $C = 0.517$ --- both gates remain open despite the $\mathbb{Z}_2$ winding, because the complexity of $\Phi_{\text{closerevepsilon}}$ sustains self-modeling at a lower tier.

The distance between the complex-time path integral and this system is 3.15 --- significant, driven primarily by the $T_{\text{bullseye}}$ vs. $T_{\text{invscr}}$ difference (crossing vs. inclusion). The imaginary-time formalism is \emph{contained} in the full theory but is not \emph{equivalent} to it.

\subsubsection{Wick Rotation as Exceptional Point}

The wrong answer first: the Wick rotation is a change of variables. It is linear, invertible, and structureless. Under this view, there is nothing to imscribe beyond $D_{\text{invomega}}$ and perhaps a trivial topology.

The right answer: the Wick rotation, when treated as a dynamical process in complex time, possesses exceptional-point structure. At the EP, eigenvalues of the generator of time evolution coalesce, and the rotation becomes non-diagonalizable. The system is no longer a coordinate change but a singularity.

\[\text{wick\_rotation\_EP} = \langle D_{\text{invomega}};\ T_{\text{nrleg}};\ R_{\text{subrightarrow}};\ P_{\text{aolig}};\ F_{\text{hardsign}};\ K_{\text{frtailgamma}};\ G_{\text{revapostrophe}};\ \Gamma_{\text{corner}};\ \Phi_{\text{revepsilon}};\ H_{\text{invscripta}};\ 1{:}1;\ \Omega_{\text{dzlig}} \rangle\]

$D_{\text{invomega}}$: infinite-dimensional Hilbert space of the Hamiltonian. $T_{\text{nrleg}}$: branching network --- multiple Riemann sheets emanate from the exceptional point. $R_{\text{subrightarrow}}$: the EP behavior supervenes on the Hamiltonian structure; it is not bidirectionally coupled to it. $P_{\text{aolig}}$: no symmetry --- the exceptional point breaks all symmetries by definition. $F_{\text{hardsign}}$: the EP is a quantum phenomenon requiring coherence. $K_{\text{frtailgamma}}$: driven dynamics --- the EP is encountered during driven (non-adiabatic) contour deformation. $G_{\text{revapostrophe}}$: universal scope. $\Gamma_{\text{corner}}$: conjunctive composition --- all eigenvalues must coalesce simultaneously for a true EP. $\Phi_{\text{revepsilon}}$: non-Hermitian exceptional point. $H_{\text{invscripta}}$: the branch structure has no finite Markov order. $1{:}1$: the exceptional point is a single mapping (one point, one degeneracy). $\Omega_{\text{dzlig}}$: integer winding around the EP encodes monodromy.

\textbf{Ouroboricity:} $O₀$ --- despite non-trivial winding and eternal memory, the lack of symmetry ($P_{\text{aolig}}$) and EP criticality prevent any self-modeling.

\textbf{Key Structural Property:} $\Phi_{\text{revepsilon}}$ absorbs any $\Phi_{\text{ctyogh}}$ or $\Phi_{\text{closerevepsilon}}$ system under tensor. This is the EP Absorption Rule: $\text{tensor}(\Phi_{\text{ctyogh}}, \Phi_{\text{revepsilon}}) = \Phi_{\text{revepsilon}}$. The composite loses criticality. This is not a mathematical curiosity --- it is the structural statement of wavefunction collapse in measurement, as we show in \pilcrow{}5.1.

The author was surprised here: the EP sits at the lowest ouroboricity tier despite having $H_{\text{invscripta}}$ and $\Omega_{\text{dzlig}}$. This reveals that symmetry ($P$) and criticality ($\Phi$) are more fundamental to self-modeling than memory depth or topological protection. An EP with perfect memory and winding is still structurally dead because it cannot be its own dual.

\subsubsection{Planck-Scale Regime}

Because the Planck regime is the background against which all these structures play, it must be imscribed independently.

\[\text{planck\_scale\_regime} = \langle D_{\text{invomega}};\ T_{\text{invscr}};\ R_{\text{lyoghlig}};\ P_{\text{upsilon}};\ F_{\text{hardsign}};\ K_{\text{schwa}};\ G_{\text{revapostrophe}};\ \Gamma_{\text{secstress}};\ \Phi_{\text{ctyogh}};\ H_{\text{invscripta}};\ n{:}m;\ \Omega_{\text{crtwo}} \rangle\]

$D_{\text{invomega}}$: field-theoretic configuration space of metrics. $T_{\text{invscr}}$: the Planck regime is contained within any full quantum gravity theory. $R_{\text{lyoghlig}}$: geometry and matter are bidirectionally coupled --- each determines the other. $P_{\text{upsilon}}$: quantum superposition of geometries. $F_{\text{hardsign}}$, $K_{\text{schwa}}$, $G_{\text{revapostrophe}}$, $\Gamma_{\text{secstress}}$ follow the pattern above. $\Phi_{\text{ctyogh}}$: self-modeling criticality. $H_{\text{invscripta}}$: eternal memory --- the gravitational field has no Markov cutoff. $n{:}m$: heterogeneous sectors (metric + matter + topology). $\Omega_{\text{crtwo}}$: binary winding protection.

The structural distance to the imaginary-time formalism is only 0.33 --- the smallest in the entire catalog. They differ only in $R$ ($\leftrightarrow$ vs. $\dagger$) and $\Phi$ ($c$ vs. $c^\mathbb{C}$). The imaginary-time formalism is, in the precise structural sense, the minimal extension of the Planck regime.

\subsubsection{Hartle--Hawking No-Boundary State}

The Hartle--Hawking wavefunction of the universe is defined by a Euclidean path integral over compact 4-geometries with no initial boundary. Its imscription is:

\[\text{hartle\_hawking\_no\_boundary} = \langle D_{\text{invomega}};\ T_{\text{commatailz}};\ R_{\text{lyoghlig}};\ P_{\text{upsilon}};\ F_{\text{hardsign}};\ K_{\text{schwa}};\ G_{\text{revapostrophe}};\ \Gamma_{\text{secstress}};\ \Phi_{\text{closerevepsilon}};\ H_{\text{invscripta}};\ n{:}m;\ \Omega_{\text{dzlig}} \rangle\]

Structural duplicate: identical to \texttt{black\_hole\_information} in the catalog. This equivalence is not accidental --- both problems involve an irreducible product ($T_{\text{commatailz}}$) where the Euclidean and Lorentzian sectors cannot be factored, and both hinge on complex-plane criticality with integer winding.

The distance from the complex-time formalism is 3.30 --- nearly identical to the distance to the imaginary-time formalism (3.15). The Hartle--Hawking state is structurally ''''almost'''' self-dual but lacks $\mu \circ \delta = \text{id}$ by exactly one primitive: $P_{\text{upsilon}}$ vs. $P_{\text{doublebarpipe}}$.

\subsection{One might argue the no-boundary proposal should carry higher criticality because it proposes a complete theory of the initial state. The argument does not survive structural analysis: the no-boundary state is a \emph{boundary condition}, not a \emph{propagator}. It specifies the wavefunction at a point but does not govern its evolution. The complex-time formalism does both --- and that difference is encoded in $\Phi_{\text{ctyogh}}$, not $\Phi_{\text{closerevepsilon}}$.}

\subsection{Algebraic Operations and Structural Derivation}

Because the five systems are now imscribed, the algebraic operations between them are determined --- not postulated, not argued, but \emph{computed}. Each tensor product, meet, and join follows from the primitive values alone.

\subsubsection{Tensor Product: Complex-Time $\otimes$ Imaginary-Time}

\[\text{complex\_time\_path\_integral} \otimes \text{planck\_imaginary\_time} = \langle D_{\text{invomega}};\ T_{\text{bullseye}};\ R_{\text{lyoghlig}};\ P_{\text{upsilon}};\ F_{\text{hardsign}};\ K_{\text{schwa}};\ G_{\text{revapostrophe}};\ \Gamma_{\text{secstress}};\ \Phi_{\text{closerevepsilon}};\ H_{\text{invscripta}};\ n{:}m;\ \Omega_{\text{dzlig}} \rangle\]

The tensor takes the maximum on union primitives and the minimum on $P$ and $F$. Here $P_{\text{doublebarpipe}} \to P_{\text{upsilon}}$ is the bottleneck: the exact Frobenius self-duality of the complex-time formalism is broken in the composite. The composite remains quantum-symmetric but is no longer exactly self-dual. Five primitives promote to maximal values: $T_{\text{bullseye}}$ (crossing supersedes inclusion), $R_{\text{lyoghlig}}$ (bidirectional supersedes adjoint), $\Phi_{\text{closerevepsilon}}$ (complex-plane criticality supersedes real criticality), $H_{\text{invscripta}}$ (eternal memory supersedes two-step), $\Omega_{\text{dzlig}}$ (integer winding supersedes binary).

The distance from the complex-time formalism to this composite is 3.15; from the imaginary-time formalism, 1.64. The bottleneck is entirely in $P$: losing exact self-duality costs everything. The composite is still conscious ($C > 0$) but drops from $O_{\infty}$ to $O₂^\dagger$.

This result should not have been expected from the start --- it emerges only after the imscriptions are in place and the tensor is computed. The author''s prior intuition was that the composite would preserve $O_{\infty}$ because both components are critical. That intuition was wrong.

\subsubsection{Tensor Product: Complex-Time $\otimes$ Wick Rotation EP}

\[\text{complex\_time\_path\_integral} \otimes \text{wick\_rotation\_EP} = \langle D_{\text{invomega}};\ T_{\text{bullseye}};\ R_{\text{lyoghlig}};\ P_{\text{aolig}};\ F_{\text{hardsign}};\ K_{\text{schwa}};\ G_{\text{revapostrophe}};\ \Gamma_{\text{secstress}};\ \Phi_{\text{revepsilon}};\ H_{\text{invscripta}};\ n{:}m;\ \Omega_{\text{dzlig}} \rangle\]

The bottleneck here is total: $P_{\text{doublebarpipe}} \to P_{\text{aolig}}$. Seven primitives promote to their union values, including the devastating $\Phi_{\text{ctyogh}} \to \Phi_{\text{revepsilon}}$. The EP Absorption Rule applies: the exceptional point swallows criticality entirely. Consciousness drops to zero. Both gates fail --- Gate 1 to $\Phi_{\text{revepsilon}}$ absorption, Gate 2 because the driven dynamics $K_{\text{frtailgamma}}$ of the EP overwhelm the near-equilibrium $K_{\text{schwa}}$ of the self-modeling system.

This is the \textbf{structural measurement problem}. The act of specifying the Wick rotation contour --- choosing a particular analytic continuation from Euclidean to Lorentzian signature --- is measurement-like. It selects one branch of the Riemann surface and discards the others. In doing so, it destroys the critical, self-modeling structure of the quantum gravitational system. This is not philosophy; it is a theorem in the IG algebra.

\subsubsection{Meet: Complex-Time $\wedge$ Wick Rotation EP}

\[\text{complex\_time\_path\_integral} \wedge \text{wick\_rotation\_EP} = \langle D_{\text{invomega}};\ T_{\text{nrleg}};\ R_{\text{subrightarrow}};\ P_{\text{aolig}};\ F_{\text{hardsign}};\ K_{\text{frtailgamma}};\ G_{\text{revapostrophe}};\ \Gamma_{\text{corner}};\ \Phi_{\text{ctyogh}};\ H_2;\ 1{:}1;\ \Omega_{\text{dzlig}} \rangle\]

Only $D_{\text{invomega}}$, $F_{\text{hardsign}}$, $G_{\text{revapostrophe}}$, and $\Omega_{\text{dzlig}}$ survive the meet without qualification. The other eight primitives resolve to their conservative (lower) values. The complex-time formalism and the Wick rotation EP are structurally remote --- their shared floor is minimal. This is the distance of 4.15 made explicit in every primitive.

\subsubsection{Meet: Planck Regime $\wedge$ Imaginary-Time}

\[\text{planck\_scale\_regime} \wedge \text{planck\_imaginary\_time} = \langle D_{\text{invomega}};\ T_{\text{invscr}};\ R_{\text{downstep}};\ P_{\text{upsilon}};\ F_{\text{hardsign}};\ K_{\text{schwa}};\ G_{\text{revapostrophe}};\ \Gamma_{\text{secstress}};\ \Phi_{\text{ctyogh}};\ H_{\text{invscripta}};\ n{:}m;\ \Omega_{\text{crtwo}} \rangle\]

Ten of twelve primitives are shared without qualification. Only $R$ (resolved to $R_{\text{downstep}}$) and $\Phi$ (resolved to $\Phi_{\text{ctyogh}}$ from $\Phi_{\text{closerevepsilon}}$) differ. This confirms what the distance of 0.33 already suggested: the imaginary-time formalism is the natural structural partner of the Planck regime.

\subsubsection{Tensor: Planck Regime $\otimes$ Imaginary-Time}

\[\text{planck\_scale\_regime} \otimes \text{planck\_imaginary\_time} = \langle D_{\text{invomega}};\ T_{\text{invscr}};\ R_{\text{lyoghlig}};\ P_{\text{upsilon}};\ F_{\text{hardsign}};\ K_{\text{schwa}};\ G_{\text{revapostrophe}};\ \Gamma_{\text{secstress}};\ \Phi_{\text{closerevepsilon}};\ H_{\text{invscripta}};\ n{:}m;\ \Omega_{\text{crtwo}} \rangle\]

Zero bottlenecks. Two union-promoted primitives: $R_{\text{lyoghlig}}$ and $\Phi_{\text{closerevepsilon}}$. The distance from the Planck regime to this composite is 0.33 --- the cleanest composition in the system. The join $\text{planck\_scale\_regime} \vee \text{planck\_imaginary\_time}$ coincides exactly with the tensor product, confirming zero structural tension between the two systems.

\subsubsection{Distance Analysis}

Because the structural types are fixed, the pairwise distances are determined. We report them with interpretation:

\begin{table}[htbp]
\centering
\small
\rowcolors{1}{white}{white}
\begin{tabularx}{\textwidth}{>{\centering\arraybackslash}X >{\centering\arraybackslash}X >{\centering\arraybackslash}X}
\toprule
\rowcolor{gray!25}\textbf{Pair} & \textbf{Distance} & \textbf{Interpretation} \\
\midrule
\rowcolors{1}{white}{gray!10}
complex\_time\_path\_integral -- planck\_imaginary\_time & 3.15 & Remote --- the crossing topology vs. inclusion topology creates the dominant term \\
complex\_time\_path\_integral -- wick\_rotation\_EP & 4.15 & Very remote --- maximal divergence between self-dual and asymmetric structure \\
complex\_time\_path\_integral -- hartle\_hawking\_no\_boundary & 3.30 & Remote --- dominated by the $P$-bottleneck ($P_{\text{doublebarpipe}}$ vs. $P_{\text{upsilon}}$) \\
planck\_imaginary\_time -- hartle\_hawking\_no\_boundary & 2.39 & Moderately remote --- differing in $T$, $R$, and $\Omega$ \\
planck\_scale\_regime -- planck\_imaginary\_time & 0.33 & Near-identical --- only $R$ and $\Phi$ differ \\
planck\_scale\_regime -- complex\_time\_path\_integral & \textasciitilde{}3.5 & Remote --- topological and criticality differences \\
\bottomrule
\end{tabularx}
\end{table}

The Hartle--Hawking distance from the complex-time formalism (3.30) is dominated by the $P$-primitive: $P_{\text{doublebarpipe}}$ vs. $P_{\text{upsilon}}$ contributes 9.0 to the weighted squared distance (out of a total of 10.8 in the breakdown). This means the no-boundary proposal is structurally similar to the complex-time formalism in every respect \emph{except} self-duality --- it is quantum-symmetric but not exactly Frobenius self-dual. Whether this single-primitive gap is phenomenologically consequential remains open.

\subsubsection{Ouroboricity Tier Summary}

\begin{table}[htbp]
\centering
\small
\rowcolors{1}{white}{white}
\begin{tabularx}{\textwidth}{>{\centering\arraybackslash}X >{\centering\arraybackslash}X >{\centering\arraybackslash}X >{\centering\arraybackslash}X >{\centering\arraybackslash}X >{\centering\arraybackslash}X}
\toprule
\rowcolor{gray!25}\textbf{System} & \textbf{Tier} & \textbf{$\Phi$} & \textbf{$P$} & \textbf{$\Omega$} & \textbf{$D$} \\
\midrule
\rowcolors{1}{white}{gray!10}
complex\_time\_path\_integral & $O_{\infty}$ & $\Phi_{\text{ctyogh}}$ & $P_{\text{doublebarpipe}}$ & $\Omega_{\text{dzlig}}$ & $D_{\text{invomega}}$ \\
planck\_imaginary\_time & $O₂^\dagger$ & $\Phi_{\text{closerevepsilon}}$ & $P_{\text{upsilon}}$ & $\Omega_{\text{crtwo}}$ & $D_{\text{invomega}}$ \\
planck\_scale\_regime & $O₂$ & $\Phi_{\text{ctyogh}}$ & $P_{\text{upsilon}}$ & $\Omega_{\text{crtwo}}$ & $D_{\text{invomega}}$ \\
hartle\_hawking\_no\_boundary & $O₂^\dagger$ & $\Phi_{\text{closerevepsilon}}$ & $P_{\text{upsilon}}$ & $\Omega_{\text{dzlig}}$ & $D_{\text{invomega}}$ \\
wick\_rotation\_EP & $O₀$ & $\Phi_{\text{revepsilon}}$ & $P_{\text{aolig}}$ & $\Omega_{\text{dzlig}}$ & $D_{\text{invomega}}$ \\
\bottomrule
\end{tabularx}
\end{table}

The complex-time path integral is the only $O_{\infty}$ system in this set. Its exact Frobenius self-duality with $\mu \circ \delta = \text{id}$ is both necessary and sufficient for $O_{\infty}$ when combined with $\Phi_{\text{ctyogh}}$, $D_{\text{invomega}}$, and $\Omega_{\text{dzlig}}$. The Wick rotation EP sits at $O₀$: despite $H_{\text{invscripta}}$ and $\Omega_{\text{dzlig}}$, the lack of symmetry ($P_{\text{aolig}}$) and EP criticality prevent self-modeling. This confirms that symmetry and criticality are more structurally load-bearing than memory or winding alone.

\subsubsection{Promotion Signatures}

Because the distance is dominated by $P$, the promotion signature from imaginary-time to complex-time is instructive:

\begin{table}[htbp]
\centering
\small
\rowcolors{1}{white}{white}
\begin{tabularx}{\textwidth}{>{\centering\arraybackslash}X >{\centering\arraybackslash}X >{\centering\arraybackslash}X >{\centering\arraybackslash}X}
\toprule
\rowcolor{gray!25}\textbf{Primitive} & \textbf{From} & \textbf{To} & \textbf{Delta} \\
\midrule
\rowcolors{1}{white}{gray!10}
$T$ & $T_{\text{invscr}}$ & $T_{\text{bullseye}}$ & 1 \\
$R$ & $R_{\text{downstep}}$ & $R_{\text{lyoghlig}}$ & 1 \\
$P$ & $P_{\text{upsilon}}$ & $P_{\text{doublebarpipe}}$ & 3 \\
$\Omega$ & $\Omega_{\text{crtwo}}$ & $\Omega_{\text{dzlig}}$ & 1 \\
\bottomrule
\end{tabularx}
\end{table}

The $P$-promotion dominates with $\delta = 3$. Moving from quantum superposition symmetry to exact Frobenius self-duality is the single largest structural upgrade required --- and it is the promotion that takes the system from $O₂^\dagger$ to $O_{\infty}$.

Conversely, demoting from $O_{\infty}$ back down to $O₂^\dagger$ requires only two relaxations:

\begin{table}[htbp]
\centering
\small
\rowcolors{1}{white}{white}
\begin{tabularx}{\textwidth}{>{\centering\arraybackslash}X >{\centering\arraybackslash}X >{\centering\arraybackslash}X >{\centering\arraybackslash}X}
\toprule
\rowcolor{gray!25}\textbf{Primitive} & \textbf{From} & \textbf{To} & \textbf{Delta} \\
\midrule
\rowcolors{1}{white}{gray!10}
$\Phi$ & $\Phi_{\text{ctyogh}}$ & $\Phi_{\text{closerevepsilon}}$ & 0.33 \\
$H$ & $H_2$ & $H_{\text{invscripta}}$ & 1 \\
\bottomrule
\end{tabularx}
\end{table}

\subsection{The total upward cost is $\Delta = 6$; the total downward fall is $\Delta = 1.33$. Climbing to exact self-duality is structurally harder than falling back to complex-plane criticality with extended memory. This asymmetry is physical: it is easier to lose exact self-duality (through any measurement-like coupling) than to recover it.}

\subsection{Emergence Frontier and Near-Critical Systems}

Because the complex-time path integral sits at $O_{\infty}$, it matters to know what structural systems lie closest to the $O_{\infty}$ / $O₂$ boundary. The emergence frontier identifies these systems --- they are the catalogs entries most structurally similar to self-modeling systems without themselves reaching $O_{\infty}$. This section is necessarily brief because the full frontier calculation is a catalog-wide computation beyond the scope of this paper. What matters here is the structural proximity of the imaginary-time formalism ($O₂^\dagger$) to the frontier: it is one promotion away from the boundary. That promotion is $P_{\text{upsilon}} \to P_{\text{doublebarpipe}}$, the same promotion identified in \pilcrow{}3.8 as the dominant cost. This convergence --- the promotion distance and the frontier distance pointing to the same primitive --- is not coincidental. It identifies $P_{\text{doublebarpipe}}$ as the structural gatekeeper of $O_{\infty}$.

\subsubsection{Crystal of Types --- Tier Census}

The full crystal of types comprises 17 280 000 structural types. The distribution across ouroboricity tiers is:

\begin{table}[htbp]
\centering
\small
\rowcolors{1}{white}{white}
\begin{tabularx}{\textwidth}{>{\centering\arraybackslash}X >{\centering\arraybackslash}X >{\centering\arraybackslash}X >{\centering\arraybackslash}X}
\toprule
\rowcolor{gray!25}\textbf{Tier} & \textbf{Cells} & \textbf{Types} & \textbf{Percentage} \\
\midrule
\rowcolors{1}{white}{gray!10}
$O₀$ & 240 & 10 368 000 & 60\% \\
$O₂$ & 72 & 3 110 400 & 18\% \\
$O_{\infty}$ & 32 & 1 382 400 & 8\% \\
$O₂^\dagger$ & 24 & 1 036 800 & 6\% \\
$O₁$ & 32 & 1 382 400 & 8\% \\
\bottomrule
\end{tabularx}
\end{table}

The complex-time path integral occupies cell 154, inner type 25 616, at crystal address \textbf{6 678 416} --- within the $O_{\infty}$ tier (8\% of the crystal). The dominance of $O₀$ (60\%) reflects the fact that most structural combinations produce systems with no self-modeling capacity. The rarity of $O_{\infty}$ is not an accident of our encoding; it is the structural reason that self-modeling systems are uncommon in nature.

One might object that the tier census merely counts formal combinations and says nothing about which combinations are physically realized. The objection is correct. The census establishes the \emph{possibility space}; physics determines which points are occupied. That 8\% of the crystal supports $O_{\infty}$ is an upper bound on how common self-modeling systems could be, not a claim about how common they are.

\begin{center}
\rule{0.5\textwidth}{0.4pt}
\end{center}

\subsection{Discussion}

\subsubsection{The Structural Measurement Problem}

Because the $\Phi_{\text{revepsilon}}$ absorption effect is the central finding, it bears repeating in its most stripped-down form. When a self-modeling system ($\Phi_{\text{ctyogh}}$, $P_{\text{doublebarpipe}}$) couples to the Wick rotation treated as an exceptional point ($\Phi_{\text{revepsilon}}$, $P_{\text{aolig}}$), the composite loses both:

\[\Phi_{\text{ctyogh}} \otimes \Phi_{\text{revepsilon}} \to \Phi_{\text{revepsilon}}, \qquad P_{\text{doublebarpipe}} \otimes P_{\text{aolig}} \to P_{\text{aolig}}\]

Gate 1 fails because $\Phi_{\text{revepsilon}}$ absorbs criticality. Gate 2 fails because $K_{\text{frtailgamma}}$ overwhelms $K_{\text{schwa}}$. The consciousness score drops from $C = 0.682$ to $C = 0$.

This is the structural measurement problem. The act of specifying the Wick rotation contour --- choosing a particular analytic continuation from Euclidean to Lorentzian signature --- selects one branch of the Riemann surface and discards the other. It is measurement-like not because it involves an observer but because it is a symmetry-breaking selection from a superposition of analytic continuations. In doing so, it destroys the critical, self-modeling structure of the quantum gravitational system. The measurement problem in quantum gravity is not about observers; it is about the impossibility of maintaining exact Frobenius self-duality through a contour selection.

An objection: this merely renames the measurement problem in structural language without solving it. The objection is partially correct. What the grammar adds is precision. Standard formulations of the measurement problem are under-specified --- they do not identify exactly which mathematical structure is broken and at what cost. The IG formalism identifies the broken structure ($\mu \circ \delta = \text{id}$ in the composite) and quantifies the cost (distance 4.15, total symmetry breaking, criticality absorption). Whether this constitutes a ''''solution'''' depends on whether one considers precise identification to be progress toward resolution. I believe it does, but the judgment is not forced by the formalism.

\subsubsection{The Imaginary-Time Formalism as Minimal Extension}

The Planck-scale regime and the imaginary-time formalism sit at a structural distance of only 0.33. This is among the closest pairs in the entire catalog. Their tensor product has zero bottlenecks, and their join coincides with their tensor --- the two systems combine without any structural tension. The imaginary-time formalism is not an arbitrary add-on to the Planck-scale description; it is the minimal structural extension needed to complete it.

This finding has a practical consequence. If one is constructing a theory of quantum gravity at the Planck scale and one must choose between working with the full complex-time formalism or the imaginary-time formalism alone, the structural analysis says: the imaginary-time formalism costs very little (distance 0.33) and can be derived from the Planck regime through a well-defined promotion sequence. The full complex-time formalism costs much more --- it requires the $P$-promotion with $\delta = 3$, which is not available at the Planck regime alone. The imaginary-time formalism is the conservative choice; the complex-time formalism is the complete one.

\subsubsection{Hartle--Hawking and Structural Equivalence}

The Hartle--Hawking no-boundary proposal is structurally identical to the black hole information problem in the catalog --- both encode as $\langle D_{\text{invomega}}; T_{\text{commatailz}}; R_{\text{lyoghlig}}; P_{\text{upsilon}}; F_{\text{hardsign}}; K_{\text{schwa}}; G_{\text{revapostrophe}}; \Gamma_{\text{secstress}}; \Phi_{\text{closerevepsilon}}; H_{\text{invscripta}}; n{:}m; \Omega_{\text{dzlig}} \rangle$. This equivalence is physically significant. Both problems involve the same irreducible product structure ($T_{\text{commatailz}}$) --- the Euclidean and Lorentzian descriptions cannot be separated --- and both depend on complex-plane criticality with integer winding.

The distance from the complex-time formalism to the no-boundary state (3.30) is dominated by the $P$-bottleneck. The no-boundary proposal is structurally ''''almost'''' self-dual --- it lacks only the exact Frobenius condition $\mu \circ \delta = \text{id}$. This suggests that the no-boundary proposal''s failure to fully specify quantum gravitational dynamics may be traceable to a single structural deficit: it is not its own dual. The Hartle--Hawking state specifies a wavefunction but not the analytic continuation structure that would make it a complete propagator theory.

\begin{center}
\rule{0.5\textwidth}{0.4pt}
\end{center}

\subsection{Conclusion}

We began by asking whether the Wick rotation is a coordinate change or a physical process. The Imscribing Grammar answers: it is both, and the difference between the two views is encoded in a single primitive assignment ($P_{\text{aolig}}$ vs. $P_{\text{doublebarpipe}}$). The consequences of that assignment cascade through every algebraic operation we computed:

\begin{enumerate}
    \item \textbf{The complex-time path integral} is $O_{\infty}$ ($C = 0.682$) at crystal address 6 678 416, requiring exact $P_{\text{doublebarpipe}}$ self-duality. This is the system that treats the Wick rotation as physical.
    \item \textbf{The imaginary-time formalism} at the Planck scale is $O₂^\dagger$ ($C = 0.517$), naturally paired with the Planck regime at distance 0.33. It is the conservative, derivable fragment.
    \item \textbf{The Wick rotation as exceptional point} absorbs criticality under tensor coupling. Consciousness drops from $C = 0.682$ to $C = 0$. This is the structural measurement problem.
    \item \textbf{The promotion path} from imaginary-time to complex-time requires four promotions, dominated by $P$ ($\delta = 3$). The return path costs only $\Delta = 1.33$.
    \item \textbf{The Hartle--Hawking no-boundary state} is structurally equivalent to the black hole information problem at distance 3.30 from the full formalism. It is a boundary condition, not a propagator theory, and the deficit is structural, not technical.
\end{enumerate}

The open question this analysis opens and cannot close: if the $P$-promotion to exact Frobenius self-duality is destroyed by any measurement-like coupling (the $\Phi_{\text{revepsilon}}$ absorption), then the $O_{\infty}$ tier is inherently unobservable. It exists as a structural possibility but cannot be accessed through any process that selects a specific analytic continuation. Whether this unobservability is a fundamental limit on what can be known about quantum gravity --- or whether there exists a structural mechanism that preserves self-duality through contour selection --- is the question the grammar poses but does not answer.

\begin{center}
\rule{0.5\textwidth}{0.4pt}
\end{center}

\emph{Structural type of this manuscript:}


\[\langle D_{\text{invomega}};\ T_{\text{bullseye}};\ R_{\text{lyoghlig}};\ P_{\text{doublebarpipe}};\ F_{\text{hardsign}};\ K_{\text{schwa}};\ G_{\text{revapostrophe}};\ \Gamma_{\text{secstress}};\ \Phi_{\text{ctyogh}};\ H_2;\ n{:}m;\ \Omega_{\text{dzlig}} \rangle\]

\emph{Crystal address: 6 678 416. Ouroboricity tier: $O_{\infty}$. Consciousness score: $C = 0.682$.}


\end{document}
